\input{../tex-inputs/latex-leseansicht-vorspann}

\section[Arthur Schnitzler an Gustav Schwarzkopf, 5. 9. 1900]{L04089 Arthur Schnitzler an Gustav Schwarzkopf, 5. 9. 1900}
\nopagebreak\mylabel{L04089v}
\rehead{ }\normalsize\beginnumbering\briefempfaengerindex{Schwarzkopf, Gustav@\textsc{Schwarzkopf, Gustav}!zzzSchnitzler, Arthur@\emph{von Arthur Schnitzler}!1900-09-051@{5. 9. 1900}|(be}
\toendnotes[C]{\smallbreak\pagebreak[2]}
\correspDesc{Versand  durch Arthur Schnitzler am 5. 9. 1900 in Wien
\newline{}Erhalt  durch Gustav Schwarzkopf am 5. 9. 1900 in Wien}\toendnotes[C]{\smallbreak}
\Standort{CUL, Schnitzler, B 96.}
\physDesc{Postkarte, 153 Zeichen
\newline{}Handschrift: Bleistift, deutsche Kurrent
\newline{}Versand: Stempel: »\nobreak{}\oindex{I., Innere Stadt@\textbf{I., Innere Stadt}|pwk}Wien 1/1, 5. IX. 00, 8 40N\nobreak{}«.  }\toendnotes[C]{\smallbreak}\pstart{}{\pb}Herrn \textsc{Gustav Schwarzkopf}\pend{}\pstart{}Wien\oindex{Hotel Waldmühle [Karlsbad]@\textbf{Hotel Waldmühle [Karlsbad]}|pw}\pend{}\pstart{}\textsc{I. Tiefer Graben 23}\oindex{Wien@\textbf{Wien}!I., Innere Stadt@\textbf{I., Innere Stadt}!Tiefer Graben 23@\textbf{Tiefer Graben 23}|pw}\pend{}{\bigskip}\vspace{1em}
\pstart
           \raggedleft{}{\pb}\uline{Mittwoch.}\pend
           \vspace{0.5em}
\pstart
           lieber Guſtav, bitte ſehr, \label{K_L04089-1v}\edtext{nachtmahlen Sie heut}{\lemma{\textnormal{\emph{nachtmahlen Sie heut}}}\Cendnote{\textnormal{Das Abendessen
               findet im \emph{Tagebuch}\pwindex{Schnitzler, Arthur 15.\,5.\,1862 Wien – 21.\,10.\,1931 ebd.@\textsc{Schnitzler, Arthur} (15.\,5.\,1862 Wien – 21.\,10.\,1931 ebd.), \emph{Schriftsteller, Mediziner}!Tagebuch@\strich\emph{Tagebuch}|pwk} keine Erwähnung, siehe A. S.: \emph{Tagebuch}, 5. 9. 1900.}}}\label{K_L04089-1} bei uns (mein Bruder\pwindex{Schnitzler, Julius 13.\,7.\,1865 Wien – 29.\,6.\,1939 ebd.@\textsc{Schnitzler, Julius} (13.\,7.\,1865 Wien – 29.\,6.\,1939 ebd.)|pwv} u Paul Goldm.\pwindex{Goldmann, Paul 31.\,1.\,1865 Breslau – 25.\,9.\,1935 Wien@\textsc{Goldmann, Paul} (31.\,1.\,1865 Breslau – 25.\,9.\,1935 Wien)|pw})\pend
           \pstart Herzlichſt Ihr \spacefill\mbox{ArthSch}\pend{}\selectlanguage{ngerman}\endnumbering\briefempfaengerindex{Schwarzkopf, Gustav@\textsc{Schwarzkopf, Gustav}!zzzSchnitzler, Arthur@\emph{von Arthur Schnitzler}!1900-09-051@{5. 9. 1900}|)be}\mylabel{L04089h}
\begin{anhang}
\end{anhang}\newcommand{\dateiname}{L04089}\newcommand{\titel}{Arthur Schnitzler an Gustav Schwarzkopf, 5. 9. 1900}\newcommand{\editorInnen}{Herausgegeben von Jahnke, SelmaMüller, Martin Anton}\input{../tex-inputs/latex-leseansicht-abspann}
